\documentclass[10pt,a4paper]{amsart}
\usepackage[margin=19mm]{geometry}
\usepackage{amsmath,amssymb,mathtools}
\usepackage[scr=rsfs,cal=euler]{mathalfa}
\usepackage{xfrac}
\usepackage[unicode,hidelinks,bookmarks=false]{hyperref}
\newcommand{\ie}{i.e.\ }
\newcommand{\cf}{cf.\ }
\newcommand{\resp}{respectively\ }
\newcommand{\Subsection}{\subsection}
\providecommand{\nobreakdash}{\nobreak\hbox{-}}
\setlength{\emergencystretch}{2em}
\multlinegap0pt
\usepackage{mathtools}
\usepackage{amssymb}
\usepackage[shortlabels]{enumitem}
\newtheorem{theorem}{Theorem}
\newtheorem{lemma}{Lemma}
\newcommand{\Z}{\mathbb Z}
\title{On zero-sum Ramsey numbers of cycles and wheels}
\author{Cheng Chi, Jialin He}
\date{Source: arXiv:2605.14954v2 (CC BY 4.0)}
\begin{document}
\maketitle

\noindent\emph{Source and licence.} On zero-sum Ramsey numbers of cycles and wheels. Version arXiv:2605.14954v2 (CC BY 4.0), distributed by arXiv under the Creative Commons Attribution 4.0 International licence, http://creativecommons.org/licenses/by/4.0/, as recorded in the arXiv licence metadata for this exact version and shown on the abstract page https://arxiv.org/abs/2605.14954v2. Full licence notice: \texttt{licenses/arxiv-2605.14954-CC-BY-4.0.txt}.\par\smallskip

\noindent\textit{Editorial scope.} Counted unit: the article's theorem thm:q3-wheel (source0.tex lines 178--180). The article's complete original proof of it is delivered with it (source0.tex lines 752--910, one \texttt{proof} environment of 159 lines, which the article places away from the statement). No mathematics is written, repaired or re-derived: the statement and the proof are the article's own lines, in the article's own notation, and no retained line is substituted or re-flowed.  Retained uncounted as the source-local context the proof needs: lines 552--560.  Nothing was omitted from any retained span, so the package carries no omission record.  No retained line contains a citation, so the wrapper carries no bibliography and no bibliography entry.  No retained line contains a figure environment, an image-inclusion command or drawing code, so no image is delivered and the package declares no asset dependency.  Editorial formatting only, in addition: an AMS \texttt{amsart} preamble that restates the article's own declarations and macros used by the retained lines, namely the environments \texttt{theorem}, \texttt{lemma} on separate counters, together with the standard packages those lines need.  The retained environments print in this document as Theorem 1, Lemma 1 in order of appearance, whereas the article prints the same items with the numbering of its own full document.  The complete native source of the article as distributed in the arXiv e-print for this exact version is bundled unchanged as \texttt{sources/200625/source0.tex}.
\begin{lemma}\label{lem:z3-affine}
    Let $w\colon E(K_N)\to \Z_3$ be an edge-labeling.  Suppose that for every ordered pair of distinct vertices $(x,y)$, the difference
    $w(xv)-w(yv)$
    is independent of $v\in V(K_N)\setminus\{x,y\}$.
    Then there exist $\lambda_v$ for every $v\in V$ and a constant $c\in\Z_3$ such that for every edge $uv$,
    \[
        w(uv)=\lambda_u+\lambda_v+c.
    \]
\end{lemma}

\begin{theorem}\label{thm:q3-wheel}
    For every $k\ge 2$, $R(W_{3k},\Z_3)=3k+1$.
\end{theorem}

\begin{proof}[Proof of Theorem~\ref{thm:q3-wheel}]
    The lower bound follows from the vertex count: $W_{3k}$ has $3k+1$ vertices.
    It remains to prove that every labeling $w\colon E(K_{3k+1})\to\Z_3$ contains a zero-sum copy of $W_{3k}$.

    Let $n=3k$ and let $V$ be the vertex set of $K_{n+1}$.
    For a vertex $h$ and a Hamilton cycle $C$ on $V\setminus\{h\}$, we define
    \[
        \Phi_h(C)=\sum_{e\in C}w(e)+\sum_{v\in V\setminus\{h\}}w(hv).
    \]
    Thus, $\Phi_h(C)$ is the weight of the wheel with center $h$ and rim $C$.  Suppose, for the contrary, that
    \begin{equation}\label{Equ:Phi ne 0}
        \Phi_h(C)\ne0
    \end{equation}
    for every $h$ and every Hamilton cycle $C$ on $V\setminus\{h\}$.

    Fix distinct vertices $x,y$.
    For $v\in V\setminus\{x,y\}$, define $\rho(v)=w(xv)-w(yv)$.
    We first claim that among the $n-1$ values $\rho(v)$, at least $n-2$ must be equal.  Indeed, otherwise, a largest value class would contain at least two vertices, and its complement would also have at least two vertices.
    We could therefore choose two disjoint pairs $\{a,c\}$ and $\{b,d\}$ such that $\rho(a)\ne\rho(c)$ and $\rho(b)\ne\rho(d)$.
    Since $n\ge6$, there exists a vertex $h\notin\{x,y,a,b,c,d\}$.
    We partition the remaining vertices of $V\setminus\{h,x,y,a,b,c,d\}$ into two ordered paths $P$ and $Q$ (which may be empty), and form a Hamilton cycle on $V\setminus\{h\}$ of the form
    \[
        C=xaPcybQdx.
    \]
    Given a path $T=t_1,\ldots,t_k$, let $T^{\rm rev}$ denote the path $t_k,\ldots,t_1$.
    Interchanging $a,c$ means replacing the subpath $aPc$ by $cP^{\rm rev}a$; this changes $\Phi_h(C)$ by $\alpha\coloneqq\rho(c)-\rho(a)\ne0$.
    Similarly, interchanging $b,d$ means replacing $bQd$ by $dQ^{\rm rev}b$; this changes $\Phi_h(C)$ by $\beta\coloneqq\rho(b)-\rho(d)\ne0$.
    By interchanging the pairs $a,c$ or $b,d$, we will get four different wheels whose weights are $S,S+\alpha,S+\beta$ and $S+\alpha+\beta$ with $\alpha,\beta\ne0$.
    In $\Z_3$ these four elements cannot all be non-zero, which contradicts~\eqref{Equ:Phi ne 0}.
    Therefore, every pair $xy$ is either regular, meaning all values of $\rho(v)$ are equal, or singular, meaning there is exactly one exceptional vertex.
    If every pair is regular, Lemma~\ref{lem:z3-affine} yields
    \[
        w(uv)=\lambda_u+\lambda_v+c.
    \]
    However, in $W_n$, every rim vertex has degree $3$, the center has degree $n\equiv0\pmod3$, and $|E(W_n)|=2n\equiv0\pmod3$.
    Hence every copy of $W_n$ has total weight zero, which is a contradiction to~\eqref{Equ:Phi ne 0}.
    Therefore, there must exist a singular pair.

    Let $xy$ be a singular pair, let $z$ be its exceptional vertex, and put $A=V\setminus\{x,y,z\}$.
    Then $|A|=n-2\ge4$.
    By the definition of a singular pair, there exist $a,b\in\Z_3$ with $a\ne b$, such that
    \begin{equation}\label{Equ:singular wheel a,b}
        w(xu)-w(yu)=a\quad(\text{for all }u\in A),
        \qquad
        w(xz)-w(yz)=b.
    \end{equation}
    Let $\delta=a-b\ne0$ and $\ell_u=w(yu)$ for $u\in A$.  For $h\in A$, define
    \[
        D_h=\sum_{v\in V\setminus\{h\}}w(hv).
    \]

    Take distinct $u,h\in A$, and let $v_1,\ldots,v_s$ be an ordering of $A\setminus\{u,h\}$, where $s=n-4$.  Consider the wheel with center $h$ and rim
    \[
        xzuyv_1\cdots v_sx.
    \]
    Interchanging $z$ and $u$ changes the wheel's weight by $a-b=\delta$.
    Since both resulting wheels have non-zero sums, the original wheel must have weight $\delta$.
    Therefore, we obtain that
    \begin{equation}\label{Equ:wheel sum}
        w(xz)+w(zu)+w(uy)+w(yv_1)+\sum_{i=1}^{s-1}w(v_i v_{i+1})+w(v_sx)+D_h=\delta.
    \end{equation}
    We define
    \begin{equation}\label{Equ:define F and H}
        F(u)=w(zu)+\ell_u \quad \text{and}\quad
        H(v_1,\ldots,v_s)=\ell_{v_1}+\sum_{i=1}^{s-1}w(v_i v_{i+1})+\ell_{v_s}.
    \end{equation}
    Substituting the above equalities into~\eqref{Equ:wheel sum}, and by~\eqref{Equ:singular wheel a,b}, there is a constant $C_0\in \Z_3$ independent of $u,h,v_1,\ldots,v_s$ such that
    \begin{equation}\label{Equ:F,D,H}
        F(u)+D_h+H(v_1,\ldots,v_s)=C_0.
    \end{equation}
    Interchanging the roles of $u$ and $h$ in~\eqref{Equ:F,D,H} yields $D_h-F(h)=D_u-F(u)$.
    Thus, we may assume that there is a constant $c_1\in\Z_3$ such that $D_i-F(i)=c_1$ for every $i\in A$.
    Therefore, \eqref{Equ:F,D,H} implies that there is a constant $C_1\in\Z_3$ such that
    \begin{equation}\label{Equ:H+F+F=C}
        H(A\setminus\{i,j\})+F(i)+F(j)=C_1
    \end{equation}
    for every pair $i,j\in A$, and the value of $H$ is independent of the ordering of its arguments.

    First, assume that $k\ge3$.
    Then $|A|\ge7$.
    Since $|A|\ge7$, for any distinct $r,p,t\in A$ we may choose $i,j\in A\setminus\{r,p,t\}$.  Applying~\eqref{Equ:H+F+F=C} to the same pair $i,j$ and to two orderings of $A\setminus\{i,j\}$ beginning with $r,p,t,\ldots$ and $p,r,t,\ldots$, respectively, using~\eqref{Equ:define F and H}, we have
    \[
        \ell_r+w(rp)+w(pt)=\ell_p+w(pr)+w(rt),
    \]
    and hence, for a fixed $t$, the value $w(rt)-\ell_r$ is independent of $r\ne t$.
    By symmetry, $w(rt)-\ell_t$ is independent of $t\ne r$.
    Since $r$ and $t$ are arbitrarily chosen from $A$, there is a constant $c\in\Z_3$ such that $w(ij)=\ell_i+\ell_j+c$ with $i,j\in A$ and $i\ne j$.
    Substituting this into~\eqref{Equ:H+F+F=C} gives
    \begin{equation*}
        2\sum_{v\in A\setminus\{i,j\}}\ell_v+c(s-1)+w(zi)+\ell_i+w(zj)+\ell_j=C_1,
    \end{equation*}
    which implies that
    \begin{equation*}
        2\sum_{v\in A}\ell_v+c(s-1)+\bigl(w(zi)-\ell_i\bigr)+
        \bigl(w(zj)-\ell_j\bigr)=C_1.
    \end{equation*}
    Thus, there exists a constant $C_2\in \Z_3$, independent of the pair $i,j$, such that
    \[
        \bigl(w(zi)-\ell_i\bigr)+\bigl(w(zj)-\ell_j\bigr)=C_2
        \qquad(\text{for all }i\ne j\in A),
    \]
    and therefore, for some constant $g\in\Z_3$, we have
    \[
        w(zi)=\ell_i+g \qquad (\text{for all }i\in A).
    \]

    Now assume that $k=2$.
    Then, $|A|=4$.
    Equation~\eqref{Equ:H+F+F=C} states that, for every partition $A=\{r,s\}\cup\{u,h\}$, there exists a constant $C_3\in \Z_3$, independent of the partition, such that $w(rs)+w(zu)+w(zh)=C_3$.
    Let $m_u=w(zu)$, this gives $w(rs)=m_r+m_s+c'$ for some constant $c'$.
    For $u,v\in A$, compare the five differences from $u$ and $v$ to the other vertices.
    By~\eqref{Equ:singular wheel a,b}, the two differences to $x,y$ are $\ell_u-\ell_v$, while the three differences to $z$ and to the other two vertices of $A$ are $m_u-m_v$.
    Since at least four of the five differences are equal, $m_u-m_v=\ell_u-\ell_v$.
    Hence, for suitable constants $c,g\in\Z_3$, for all $u\in A$, $m_u=\ell_u+g$ and again
    \[
        w(uv)=\ell_u+\ell_v+c,
        \qquad
        w(zu)=\ell_u+g
        \qquad(\text{for all } u,v\in A,\ u\ne v)
    \]


    In all cases, for $u,v\in A$, $u\ne v$, we have
    \begin{equation}\label{Equ:wheel structure}
        w(yu)=\ell_u,
        \qquad
        w(xu)=\ell_u+a,
        \qquad
        w(zu)=\ell_u+g,
        \qquad
        w(uv)=\ell_u+\ell_v+c.
    \end{equation}
    Define a weight $q_0(uv)=\lambda_u+\lambda_v+c$ by
    \[
        \lambda_u=\ell_u\ (\text{for all }u\in A),
        \qquad
        \lambda_y=-c,
        \qquad
        \lambda_x=a-c,
        \qquad
        \lambda_z=g-c.
    \]
    Then, $q_0$ agrees with $w$ on every edge incident with at least one vertex of $A$.  Replacing $w$ with $w-q_0$ preserves the weight of every $W_n$, because every vertex degree in $W_n$ is divisible by $3$ and $|E(W_n)|=2n\equiv0\pmod3$.
    Hence, we may assume that all edges incident with $A$ have weight $0$.
    Only the three edges $xy,xz,yz$ may be non-zero.

    We choose a center $h\in A$.
    Since $|A\setminus\{h\}|=n-3\ge3$, we can order the rim vertices so that $x,y,z$ are pairwise non-adjacent.
    For instance, we may choose $xa_1ya_2za_3\cdots x$
    with $a_1,a_2,a_3\in A\setminus\{h\}$ distinct and all remaining vertices of $A\setminus\{h,a_1,a_2,a_3\}$ inserted in the final path.
    This wheel uses no edge among $x,y,z$, and all its other edges are incident with $A$.
    Hence, its total weight is zero, contradicting the assumption~\eqref{Equ:Phi ne 0}.

    Therefore, every labeling of $K_{3k+1}$ contains a zero-sum $W_{3k}$, and
    \[
        R(W_{3k},\Z_3)=3k+1,
    \]
    which completes the proof of Theorem~\ref{thm:q3-wheel}.
\end{proof}

\end{document}
